\section{Méthodes de classification}

\begin{frame}{Approche double}

\begin{columns}
\begin{column}{0.48\textwidth}
\textbf{Clustering classique}

\vspace{1em}
\begin{center}
\begin{tikzpicture}[
  box/.style={draw, rounded corners=4pt, minimum width=2.2cm, minimum height=0.7cm,
              align=center, font=\tiny, thick},
  arr/.style={-{Stealth[length=3pt]}, thick, color=pipeGray},
  scale=0.7, transform shape
]
  \node[box, fill=pipeBlue!10, draw=pipeBlue!50] (feat) at (0, 0)
    {\textbf{Matrice 2516×25}\\moy |dG|};

  \node[box, fill=pipeOrange!10, draw=pipeOrange!50] (cos) at (3, 0)
    {\textbf{Cosine similarity}\\k-NN graph};

  \node[box, fill=pipeGreen!10, draw=pipeGreen!50] (clust) at (6, 0)
    {\textbf{Leiden / Spectral / GN}\\classification};

  \draw[arr] (feat) -- node[above, font=\tiny] {} (cos);
  \draw[arr] (cos) -- node[above, font=\tiny] {} (clust);
\end{tikzpicture}
\end{center}

\vspace{0.5em}
\small
Clustering spectral (k fixe) et Girvan-Newman (edge-betweenness).
\end{column}

\begin{column}{0.48\textwidth}
\textbf{Décomposition matricielle (NMF)}

\vspace{1em}
\begin{center}
\begin{tikzpicture}[
  box/.style={draw, rounded corners=4pt, minimum width=2.2cm, minimum height=0.7cm,
              align=center, font=\tiny, thick},
  arr/.style={-{Stealth[length=3pt]}, thick, color=pipeGray},
  scale=0.7, transform shape
]
  \node[box, fill=pipeBlue!10, draw=pipeBlue!50] (feat) at (0, 0)
    {\textbf{Matrice 2516×25}\\moy |dG|};

  \node[box, fill=pipeOrange!10, draw=pipeOrange!50] (res) at (3, 0)
    {\textbf{Résidus}\\V - col\_mean, clip+};

  \node[box, fill=pipePurple!10, draw=pipePurple!50] (nmf) at (6, 0)
    {\textbf{NMF}\\$V \approx WH$};

  \draw[arr] (feat) -- node[above, font=\tiny] {} (res);
  \draw[arr] (res) -- node[above, font=\tiny] {} (nmf);
\end{tikzpicture}
\end{center}

\vspace{0.5em}
\small
Soft membership : chaque miRNA a un score par signature.
\end{column}
\end{columns}
\end{frame}

\begin{frame}{Grid search}

\begin{block}{Paramètres testés}
\begin{itemize}
  \item \textbf{Power} : 1 (linéaire), 2 (quadratique), 3 (cubique) — poids = $\lvert$dG$\rvert^p$
  \item \textbf{Dimensions UMAP} : 2, 3, 4, 5
  \item \textbf{Résolution Leiden} : balayée pour maximiser silhouette + modularité
\end{itemize}
\end{block}

\vspace{0.5em}
\textbf{Résultat optimal} : power = 1, dimensions = 2 (résultats stables, meilleur ratio qualité/complexité).
\end{frame}